
\question
Many ministers believe that praying for sick people will speed their recovery.  This is

\begin{choices}
\choice a scientific hypothesis because sick people get better every day and people have prayed for them.
\choice not a scientific hypothesis because the presence of a God that could answer prayer can't be tested with an experiment.
\choice not a scientific hypothesis because prayer is a supernatural phenomenon that can't be tested with an experiment.
\correctchoice a scientific hypothesis because the statement makes a specific claim that can be tested with an experiment.
\choice a scientific hypothesis because praying is a natural process that people do every day.
\end{choices}

\question
When you perform an experiment, the results are compared to the:

\begin{choices}
\choice Observation 
\choice Conclusion
\choice Experiment
\choice Hypothesis
\correctchoice Prediction
\end{choices}

\question
Which of the following most accurately describes the process by which scientific knowledge advances? 

\begin{choices}
\choice The use of inductive reasoning to explain specific experimental results about the natural world.
\choice The use of scientific analysis to explain the natural world by theory and experiments.
\correctchoice The objective analysis of the natural world by hypothesis testing and experiments.
\choice The analysis of the natural world through theoretical reasoning and experiments.
\choice The use of deductive reasoning to develop broad scientific principles that explain the natural world.
\end{choices}

\question
If female birds are given a limited amount of food, then they will produce fewer eggs.  This is an example of

\begin{choices}
\choice an observation that female birds produce few eggs when food is scarce.
\correctchoice a prediction of the experimental results that will support the tentative explanation for why female birds produce fewer eggs when food is scarce.
\choice a hypothesis because it is a tentative explanation for why female birds produce fewer eggs when food is scarce.
\choice a hypothesis that explains the correlation between number of eggs produced and food availability.
\choice a fact that the bird's diet includes compounds necessary to produce eggs.
\end{choices}

\question
Which of the following is the best definition of a hypothesis?

\begin{choices}
\correctchoice A tentative explanation for a natural phenomenon.
\choice A tentative explanation about a supernatural phenomenon.
\choice A prediction about the expected results from a specific experiment.
\choice A brief summary of a broad scientific theory.
\choice An alternative explanation for a null hypothesis.
\end{choices}

\question
Prematurely born babies tend to have a higher death rate because they have not received all of the nutrients they need for proper development. This is an example of

\begin{choices}
\correctchoice a hypothesis because it is a tentative explanation for high death rate of prematurely born babies.
\choice a fact that nutrition is important for fetal development.
\choice an observation that prematurely born babies experience high death rates.
\choice an experiment that tests for death rate in prematurely born babies.
\choice a theory that explains the well-known phenomenon of high death rate in prematurely born babies because it has been tested many times and has always been supported by the evidence.
\end{choices}

\question
What is the importance of randomization when designing a scientific experiment?

\begin{choices}
\choice Randomization is required to test between the null and alternative hypotheses.
\choice Randomization is required to test between treatment and control groups.
\correctchoice Randomization is required to avoid introducing unwanted bias into the experimental results.
\choice Randomization is required to determine the proper significance level to use, such as $p \le 0.05$.
\choice Randomization is required to identify the independent and dependent variables. Independent variables are not randomized. Dependent variables are randomized.
\end{choices}

\question
The Theory of Natural Selection, a basic theory of evolutionary biology, states that certain traits become more common in a population because the traits allow some individuals to produce more offspring. The Theory of Natural Selection is a theory because

\begin{choices}
\correctchoice it has been tested many times with many different hypotheses and has always been supported.
\choice it is not a fact, and may or may not be true.
\choice it is a tentative explanation of about how organisms adapt to their environment.
\choice it has been tested many times with many different hypotheses and has been proved to be true.
\choice evolution can't be observed, so it can't be tested. It is basically a hypothesis.
\end{choices}

\question
In science, a hypothesis is only useful if it

\begin{choices}
\correctchoice can be tested.
\choice explains the natural world.
\choice agrees with the conclusion.
\choice can be proven correct.
\choice is proven correct.
\end{choices}

\question
When Abby Hodgins reported her experimental results, she included the probability value  $p < 0.05$.  This value indicates the probability that her results were due to her experiment treatment.

\begin{choices}
\choice True%dontshuffle
\correctchoice False%dontshuffle
\end{choices}

\newpage

\question
The figure shows the amount of nitrogen contained in water runoff in a deforested area compared to a non-deforested, control area. Which of the following is the correct interpretation of this graph?
\includegraphics[width=0.9\textwidth]{runoff}

\begin{choices}
\choice When trees are cut, they produce high amounts of nitrates that end up in the runoff.
\correctchoice Deforestation increases the amount of runoff that contains nitrates.
\choice The control experiment decreased the nitrogen concentration in the runoff.
\choice Tree cutting does not affect nitrate concentration because the deforested and control areas have the same concentration soon after tree cutting was complete.
\choice The nitrate concentration in runoff increased by as much as $8\times$ higher in deforested areas compared to the control.
\end{choices}

\question
What is the probability of flipping a coin three times and getting three tails? \textbf{Show your work. 3 points.}

\begin{choices}
\choice $1/16$
\choice $1/6$
\choice $1/4$
\correctchoice $1/8$
\choice $1/2$
\end{choices}

\question
Which of the following questions cannot be answered by the scientific method?

\begin{choices}
\choice Can rats vomit?
\choice Which fuel source is the most economical to replace coal?
\correctchoice Who deserves to be vaccinated in a time of crisis?
\choice What is the most nutritious food for cats?
\choice How far is the planet Neptune from the Sun?
\end{choices}

\newpage

\question
After invention of the microscope, scientists discovered that all living organisms were composed of one or more cells, that the cell was the basic unit of function in all organisms, and that only existing cells gave rise to new cells. These specific observations were eventually organized into a general theory of biology called the Cell Theory. This is an example of 

\begin{choices}
\choice deductive reasoning.
\choice objective reasoning.
\correctchoice inductive reasoning.
\choice theoretical reasoning.
\choice scientific reasoning.
\choice hypothetical reasoning.
\end{choices}


\question
You roll a pair of dice.  What is the probability of getting a 2 and a 5, in any order? \textbf{Show your work.  3 points.}

\begin{choices}
\choice $1/3$
\correctchoice $1/18$
\choice $1/10$
\choice $1/2$
\choice $1/36$
\end{choices}

\question
Which of the following statements \emph{cannot} be tested by the scientific method?

\begin{choices}
\correctchoice God created light and dark, then stars.
\choice The order in which organisms, including humans, first appeared on earth. For example, man, then plants, then land animals and birds, then woman.
\choice All of these statements are religious claims so none can be tested by the scientific method.
\choice The Earth is between 6000 and 10,000 years old.
\choice The entire universe, the earth, and all living things were created in six days.
\choice The Grand Canyon was formed by a global flood about 4000 years ago.
\end{choices}

\question
A scientist predicted that fertilizer added to soil would cause plants to produce more fruit compared to soil that did not get fertilizer. After her experiment, plants that received fertilizer did produce more fruit than the other plants. Did she \emph{prove} that fertilizer always causes plants to produce more fruit? (Choose the best answer.)

\begin{choices}
\choice Maybe. She should do more experiments to be sure.
\choice Yes. The results of her experiment showed that adding fertilizer always causes plants to grow more fruit.
\correctchoice No. She showed that fertilizer caused those particular plants to produce more fruit. There is always a chance that other plants may not produce more fruit when given fertilizer.
\choice Yes, because her hypothesis predicted that the plants would produce more fruit if fertilizer was added to the soil.
\choice No. Her hypothesis was falsified because fertilizer causes plants to grow taller but doesn't affect fruit production.
\end{choices}

\question
Why must an experiment include both treatment and control groups?

\begin{choices}
\choice To determine the standard deviation around the average.
\correctchoice To determine whether the results are significantly different from natural random variation.
\choice To include both independent and dependent variables in the experiment.
\choice To test both the null and alternative hypotheses.
\choice To avoid introducing unwanted bias into the experiment.
\end{choices}

\question
Standard deviation provides a measure of the amount of variation around the average value. About how much of the variation is contained within 1 standard deviation above and below ($\pm$) the average?

\begin{choices}
\choice 58\%
\choice 78\%
\choice 65\%
\choice 55\%
\correctchoice 68\%
\end{choices}

